\chapter{引言}

\section{研究背景和意义}
随着高校的教学规模不断扩大与教学平台的数字化进程不断普及，课程教学不再局限于传统课堂讲授的边界，广泛延伸到线上讨论和作业系统等多元场景\cite{zawacki2019systematic}。学生在预习、复习、完成作业和巩固知识点时会自然产生大量个性化疑问，而教师和助教受到时间和精力限制，难以逐一满足所有同学的求助需求。尤其是在程序设计、数据结构等实践性较强的课程中，学生需要厘清概念性知识点，还得解决代码调试、实验步骤梳理、实验环境配置等具体问题。传统的以人工为主的答疑方式在响应时效和回复一致性上存在固有瓶颈，拉低学生的学习体验，也会增加教师的教学负担。

我国正在持续推进教育数字化建设，教育部自2022年起实施国家教育数字化战略行动，强调通过数字技术扩大优质教育资源供给，提升教育服务能力和教学管理效率，教育系统的信息化建设正逐步转向以数据、平台和智能服务为支撑的数字化转型\cite{2022教育部2022年工作要点部署实施教育数字化战略行动}。课程助教系统正是在这样的背景下产生的，它可以在一定程度上缓解教师重复答疑压力并为学生提供更加及时的学习支持。

大语言模型（Large Language Model, LLM）在自然语言理解、文本生成和对话交互上的能力跃升使其能够以更加贴近人类思维的方式理解学生的问题并生成解释性回答，为构建面向课程的智能助教系统提供技术基础。然而单纯依赖通用大语言模型难以满足课程答疑的要求，它给出的答案很容易脱离实际的课程内容，生成一些表面合理却与事实不符的内容，即“幻觉”问题\cite{ji2023survey}。检索增强生成（Retrieval-Augmented Generation, RAG）\cite{lewis2020retrieval} 不单纯依赖模型内部知识，将教材文本、课程资料与实验手册等整合为外部知识库，在生成回答时动态检索并引用相关片段，从而将答案约束在课程内容的边界内。这一做法让回答的可靠性建立在可直接追溯的课程原文之上，“幻觉”问题也因“检索”得到有效抑制。将 LLM 与课程知识库深度融合，系统能够产出与真实教学内容高度一致的解释，为构建面向课程场景的智能助教系统提供了一条更严谨、更具可溯性的技术路径。

针对上述教学场景下的答疑供需矛盾与通用大语言模型在课程场景中的幻觉风险，构建一个基于 RAG 与知识溯源技术的 AI 助教系统具有明确的理论价值与实践意义，让学生在学习过程中随时获得基于教材与课件的可靠回答，提升学习效率，系统提供的引用来源能够帮助学生理解知识出处，避免依赖不可靠的生成内容，可以减轻教师在重复性答疑上的负担，提高教学效率，为教学辅助工具的建设提供可行的技术方案。 

\section{国内外研究现状}
国外关于课程助教与智能教学支持的研究起步较早，早期探索主要围绕智能辅导系统（ITS）展开，Kulik 和 Fletcher 对 50 项受控研究进行元分析\cite{kulik2016effectiveness}后指出智能辅导系统相较传统教学方式能够带来较为稳定的学习增益，但这一效果高度依赖测试目标与教学目标的对齐程度，表明教学系统的价值在于回答是否真正服务于具体学习任务。

进入 LLM 阶段后，研究重心从固定规则驱动逐步转向自然语言理解、内容生成与个性化支持，Kasneci 等指出 LLM 可以生成教学内容、提升互动性并支持个性化学习，但同时也带来了偏差、幻觉、可解释性不足和人类监督缺失等风险\cite{kasneci2023chatgpt}，Xu 等进一步从教育应用角度总结了 LLM 在教学、反馈和学习支持中的潜力，同时指出其落地仍受制于可靠性、评估标准和场景适配\cite{xu2024large}。正是在这一背景下，RAG 被视为提升教育问答可信度的重要路径，Lewis 等提出的 RAG 方法通过引入外部知识检索使生成结果更具体、更事实化，也为教育场景中“有依据的回答”提供了技术基础\cite{lewis2020retrieval}。UNESCO 在 2023 年发布的指导意见则从更高层面划定了边界，强调生成式人工智能进入教育后必须坚持以人为中心，并持续关注隐私、伦理、透明性与教学适配问题\cite{holmes2023guidance}。

国内研究的推进则与教育数字化和教育大模型建设紧密交织，在国家教育数字化战略行动的持续驱动下，关注点已从资源平台建设逐步转向智能化服务与教学流程重构\cite{2022教育部2022年工作要点部署实施教育数字化战略行动}。曹培杰等指出教育大模型并非通用大模型的简单微调，而是知识库技术、智能教育技术与大模型能力的集成，其核心在于面向教育场景持续提供专用支持\cite{曹培杰2024教育大模型的发展现状、创新架构及应用展望}。余胜泉和熊莎莎强调基于大模型增强的通用人工智能教师架构如果无法解决幻觉、深层逻辑缺失和社会情感缺失等问题就很难真正进入教学过程\cite{余胜泉2024基于大模型增强的通用人工智能教师架构}。肖建力等在对智慧教育中 LLM 的综述中也指出国内研究已开始从技术可行性转向教育场景适配、定制化训练和应用边界探讨，但整体仍偏重架构讨论与概念验证，面向具体课程场景的系统实现和效果评估尚不充分\cite{肖建力2025智慧教育中的大语言模型综述}。

国外研究在智能辅导系统、教育大模型和检索增强生成方面起步较早，已沉淀出相对清晰的技术脉络，国内研究则在教育数字化政策的推进下与之呼应，更多聚焦于教育大模型、智慧教育和智能教师架构等方向。现有工作已充分验证了 LLM 进入教育场景的可能性，但几个关键问题仍未被细致拆解：如何让模型只在课程知识的边界内作答，如何把检索、提示和生成组织成一条稳定的流程，以及如何在具体课程中验证系统的实际效果。本文的课程助教系统，正是瞄准这些尚未被充分回答的问题展开的。

\section{研究内容}
本文围绕基于 RAG 和知识溯源的课程助教系统展开，处理课程资料的组织、学生问题的理解以及模型回答的约束这三个环环相扣的问题。系统的设计从课程助教的实际使用场景出发，面向知识范围明确的教学场景，给出的回答必须被限定在课程资料和教学逻辑所允许的范围内。基于这一前提，本文对学生常见提问类型、课程资料形态以及系统需承担的辅助职责进行了梳理，为后续实现提供需求依据。系统总体架构将前端交互、后端处理、知识检索和模型生成四个层次整合起来，从用户输入、资料召回到提示组织与结果输出形成一个闭环，防止模型直接面对原始问题时出现知识来源不清和回答范围失控，并通过外部知识支撑来强化回答的针对性与一致性。

围绕系统实现，本文重点推进了四项相互衔接的工作：对课程相关资料进行整理与处理，建立起适用于检索的知识库；设计检索与提示词的组合策略，让模型在生成回答前先获得与问题高度相关的上下文信息；实现课程助教的交互界面与后端调用逻辑，保证学生可以自然地用语言提问并获取反馈；通过功能测试与典型案例分析，从可用性、响应效果和回答质量三个维度验证系统的实际表现。以课程场景中的真实需求为出发点，以检索增强生成机制为技术主轴，以系统实现和效果验证为落脚点，构建一个既能服务课程教学、又始终被课程知识约束的智能助教系统。


\section{论文组织结构}
？？？改
本文共分为六章，各章内容安排如下。
第一章为绪论，主要介绍课题研究背景、研究意义、国内外研究现状、研究内容和论文结构。
第二章为相关技术基础，介绍大语言模型、检索增强生成、提示词工程、向量检索以及系统开发相关技术。
第三章为需求分析与系统总体设计，分析系统功能需求和非功能需求，并给出系统总体架构和核心业务流程。
第四章为系统详细设计与实现，重点说明知识库构建、检索模块、提示词设计、模型调用和前端交互界面的实现过程。
第五章为系统测试与效果分析，通过测试用例和实验结果分析系统功能、性能和回答质量。
第六章为总结与展望，总结本文完成的主要工作，分析系统不足，并提出后续改进方向。